\section{The two original prime actions as a finite subscheme of the norm-one torus}
\label{sec:two-prime-torus}

The norm-one torus in (AT25) has an exact receiver in the original
arithmetic zero classes. The construction below uses the actual primes
$2,3$, the original coordinate $s$, and every multiplicity. It does not
replace a prime eigenvalue by its modulus. The torus comes from the
explicit descent of the Connes--Consani twistor charts
\cite{ConnesConsaniTwistor}; the full residue and prime actions used here
are those proved in the full-jet chapter and in (PL23)--(PL28).

Let $Z$ be a nonempty finite set of actual zeros of $g=2\xi$, stable
under $\iota\rho=1-\bar\rho$, with each full multiplicity $m_\rho$.
Write
\[
 A_Z=\prod_{\rho\in Z}\C[t_\rho]/(t_\rho^{m_\rho}),
 \qquad t_\rho=s-\rho,
 \qquad f^\dagger(s)=\overline{f(1-\bar s)},
 \qquad A_{Z,\R}=\{f:f^\dagger=f\}.
 \tag{TP1}\label{tp:algebra}
\]
Every $f\in A_Z$ has the unique decomposition
$f=(f+f^\dagger)/2+i(f-f^\dagger)/(2i)$ into two elements of
$A_{Z,\R}$. Consequently $A_{Z,\R}\otimes_\R\C\simeq A_Z$,
with the given conjugation. This identity keeps the actual reflection;
it does not conjugate each off-critical factor separately.

For each original prime $p$ put
\[
 \lambda_p=p^{-1/2}\Pi_p(1)
 =\bigl(p^{\rho-1/2}\sum_{j=0}^{m_\rho-1}
       (\log p)^jt_\rho^j/j!\bigr)_{\rho\in Z},
 \quad
 a_p=\frac{\lambda_p+\lambda_p^{-1}}2,
 \quad b_p=\frac{\lambda_p-\lambda_p^{-1}}{2i}.
 \tag{TP2}\label{tp:coordinates}
\]
Here $\Pi_p$ is multiplication by the full $p^s$ jet. The factor
$p^{-1/2}$ is stated explicitly: the original action is recovered
exactly as $p^{1/2}\lambda_p$. The exponential and its inverse are
finite Taylor expressions in each original nilpotent coordinate.
Reflection gives $\lambda_p^\dagger=\lambda_p^{-1}$, whence
$a_p^\dagger=a_p$, $b_p^\dagger=b_p$, and $a_p^2+b_p^2=1$.
There is therefore a real algebra morphism
\[
 \Psi:\mathcal B:=
 \R[a_2,b_2,a_3,b_3]/(a_2^2+b_2^2-1,a_3^2+b_3^2-1)
 \longrightarrow A_{Z,\R}
 \tag{TP3}\label{tp:map}
\]
with exactly the images in TP2. Its source is the coordinate algebra
of the square of the real norm-one torus, not a freely assigned
positivity condition.

\begin{theorem}[The two primes recover every original jet]
The morphism $\Psi$ is surjective. Thus the entire real packet is
the finite closed subscheme of $\mathcal T^2$ defined by
$I_Z=\ker\Psi$. Its complexification has coordinate maps
$X_2\mapsto\lambda_2$, $X_3\mapsto\lambda_3$ from
$\C[X_2^{\pm1},X_3^{\pm1}]$. No positive-order jet is discarded.
\end{theorem}
\begin{proof}
Two distinct roots $\rho,\sigma$ cannot have equal values for both
$\lambda_2$ and $\lambda_3$. Otherwise
$(\rho-\sigma)\log2=2\pi i n$ and
$(\rho-\sigma)\log3=2\pi i m$ for integers $m,n$. A nonzero
difference would make $m,n$ nonzero and give $m\log2=n\log3$,
contradicting unique prime factorization after exponentiating
(negative exponents can be cleared). Thus the two values separate
the roots.

Choose a complex number $c$ outside the following finite set.
For every pair of distinct roots exclude a value solving
$\lambda_2(\rho)+c\lambda_3(\rho)
=\lambda_2(\sigma)+c\lambda_3(\sigma)$, if there is one.
For every root exclude the unique value solving
$(\log2)\lambda_2(\rho)+c(\log3)\lambda_3(\rho)=0$.
Then $h=\lambda_2+c\lambda_3$ has distinct constant values
$h_\rho$ and nonzero first derivative at every root. In the local
factor, $h-h_\rho=d_\rho t_\rho+O(t_\rho^2)$ with
$d_\rho\ne0$. Successively matching the coefficients of
$t_\rho,t_\rho^2,\ldots,t_\rho^{m_\rho-1}$ expresses
$t_\rho$ as a polynomial in $h-h_\rho$: at stage $j$ the
coefficient to be divided by is $d_\rho^j\ne0$. This proves
that the local map from $\C[X]$ has kernel
$(X-h_\rho)^{m_\rho}$ and is onto.

For distinct roots these ideals are coprime. Bezout's identity for
their powers gives polynomials equal to one modulo one power and
zero modulo all others. Combining these with the local polynomial
expressions proves that $\C[h]=A_Z$. In particular the image of
$\C[X_2^{\pm1},X_3^{\pm1}]$ is all of $A_Z$. Complexifying TP3
is that map, because $X_p=a_p+ib_p$, $X_p^{-1}=a_p-ib_p$.
For a real fixed element choose a complex polynomial preimage and
average it with its real conjugate. The average lies in $\mathcal B$
and still maps to that element. This proves surjectivity over $\R$.
The auxiliary $c$ proves generation; it changes neither original
prime nor either of their actions.
\end{proof}

There is also an explicit, exact description of the ideal, useful
for retaining the full multiplicities:
\[
 I_Z\otimes_\R\C=
 \left\{P:\left.
 \frac{d^j}{ds^j}P(2^{s-1/2},3^{s-1/2})
 \right|_{s=\rho}=0\quad
 (\rho\in Z,\ 0\le j<m_\rho)\right\}.
 \tag{TP4}\label{tp:ideal}
\]
These are exactly the coefficients of the original Taylor map, so
the formula follows by its kernel definition, without a radical
replacement. For any other rational prime $p$, TP2 belongs to
this same quotient. The constructive Bezout and local inversion
proof above expresses it as a polynomial in the two original
prime coordinates. Therefore $\Pi_p$ on this packet is recovered
by that polynomial followed by multiplication by $p^{1/2}$.

\begin{proposition}[The real closed points and their full thickening]
As a real algebra, the packet has the exact factorization
\[
 A_{Z,\R}\simeq
 \prod_{\rho=\iota\rho}\R[y_\rho]/(y_\rho^{m_\rho})
 \ \times\!
 \prod_{\{\rho,\iota\rho\},\,\rho\ne\iota\rho}
       (\C[t_\rho]/(t_\rho^{m_\rho}))_{\R},
 \qquad y_\rho=t_\rho/i.
 \tag{TP5}\label{tp:real-factors}
\]
The real closed points of $\Spec A_{Z,\R}$ are exactly the
critical-line roots. Each off-critical reflection pair is a
closed point with residue field $\C$ and the displayed thickening.
At a real point $|\lambda_p|=1$ for every prime. At a root with
$\Re\rho\ne1/2$ the modulus is exactly $p^{\Re\rho-1/2}$.
\end{proposition}
\begin{proof}
At a fixed root $t_\rho^\dagger=-t_\rho$, so $y_\rho$ is
fixed. A polynomial in $y_\rho$ is fixed exactly when all its
coefficients are real. At a two-element orbit, one arbitrary
complex truncated polynomial determines its partner uniquely by
conjugating coefficients and replacing $t$ by $-t$. Projection
onto the chosen factor is a real algebra isomorphism. These
computations prove TP5. The maximal ideal in each truncated
factor is its positive-order ideal, giving the claimed residue
fields. There is no unital real algebra map $\C\to\R$, since
the square of the image of $i$ would equal $-1$. Thus only the
first factors give real-valued points. Finally the modulus follows
directly from the original exponential in TP2. The real structure
on the complex torus is $\lambda\mapsto\bar\lambda^{-1}$;
its nonfixed points occur in precisely the pairs in TP5.
\end{proof}

This calculation specifies the purity condition geometrically:
it is the assertion that the relevant reduced arithmetic packet
has only real closed points in this real torus. Merely supplying
its real structure does not establish that assertion; TP5 retains
and describes every nonreal point rather than deleting it.
Its nilpotent part also remains explicit. With the original
residue pairing and Jacobian, the exact pullback is
\[
 R_Z(f,J_g h)=
 \sum_{\rho\in Z}m_\rho\overline{f(\iota\rho)}h(\rho),
 \qquad \ker J_g=\prod_\rho(t_\rho),
 \tag{TP6}\label{tp:trace}
\]
as proved in PL24--27. Thus a fixed value contributes a positive
one-dimensional block, while each two-element orbit contributes
$\left(\begin{smallmatrix}0&m_\rho\\m_\rho&0\end{smallmatrix}\right)$,
of signs $+,-$. Every positive-order jet stays in the stated
trace radical and in the perfect raw residue space. Surjectivity
in TP3 means these are the signs and jets seen by the two actual
prime-coordinate functions, not an unrelated choice of tests.

\subsection{Its exact map over the full split-support spectrum}
Let $L$ be the original finite nontrivial distributive support
lattice. Applying its coefficient operations to TP3 gives the
surjective semiring morphism
\[
 G_L(\mathcal B)\longrightarrow G_L(A_{Z,\R}),\qquad
 (b,1_L)\longmapsto(\Psi b,1_L),\quad
 z_\lambda\longmapsto z_\lambda.
 \tag{TP7}\label{tp:split}
\]
Every killed arithmetic coefficient goes to the supported zero
$e$, and every original support label is fixed. On prime spectra
the induced map is the topological embedding
\[
 \SpecLat L\triangleright\Spec A_{Z,\R}
 \ \hookrightarrow\
 \SpecLat L\triangleright\Spec\mathcal B,
 \tag{TP8}\label{tp:spectrum}
\]
equal to the identity on support primes and to the contraction
of the closed packet $V(I_Z)$ on the arithmetic part.
Indeed a support prime omits $e$, so a coefficient mapping to
$e$ cannot enter it. An arithmetic prime contains every support
label, and its supported coefficients contract by $\Psi$.
These observations give exactly TP8 using the prime classification
of the foundational split-support paper \cite{support}.
Surjectivity of TP7 shows that the basic open sets in its domain
spectrum are the restrictions of basic opens from the target,
which proves the embedding assertion.

The topology of this image differs from that of the ring packet
in a precise way: its image contains the whole support stratum,
and its closure is the whole target spectrum. Each support point
specializes to every arithmetic point in the ordinal-sum topology.
Since $L$ is finite and nontrivial, it has a support prime: choose
a minimal nonzero join-irreducible by finite descent, and use its
prime filter. Thus every arithmetic point lies in that closure,
and every support point is already in the image.
The image is proper. For example, the real torus has infinitely
many pairs of real points, while the finite-dimensional algebra
$A_{Z,\R}$ has only finitely many prime ideals by TP5. Evaluation
at those torus points gives infinitely many distinct maximal
ideals of $\mathcal B$, so they cannot all belong to $V(I_Z)$.
Consequently TP8 is a proper dense embedding, with its finite
arithmetic part closed inside $V(e)$. No ring closed-immersion
statement has been transferred to the entire split spectrum.

\begin{figure}[ht]
\centering
\begin{tikzcd}[column sep=large,row sep=large]
 \SpecLat L\triangleright\Spec A_{Z,\R}
 \arrow[r,hook,"\text{TP8}"]
 &\SpecLat L\triangleright\Spec\mathcal B\\
 \Spec A_{Z,\R}\arrow[u,hook]\arrow[r,hook,"V(I_Z)"]
 &\mathcal T^2\arrow[u,hook]
\end{tikzcd}
\caption{The original two primes determine the full finite jet
packet by TP2--5. The bottom map is a closed immersion of real
schemes. The top map retains every support prime and is a proper
dense topological embedding. Both vertical maps identify the
arithmetic parts inside $V(e)$. The change is proved directly
from the foundational ordinal-sum spectrum; the diagram does not
assign positivity to its nonreal closed points.}
\end{figure}
